\section*{Chapter \ref{ch:m2:repo-and-specials} --- Repo and Specials}

\begin{solution}{exo:m2:repo-and-specials:1}
Friday to Tuesday is four days: $50\,000\,000 \times 0.039 \times 4/360 =
\text{USD}~21\,666.67$.
\end{solution}

\begin{solution}{exo:m2:repo-and-specials:2}
Collateral value $100\,000\,000 \times 1.015 = 101\,500\,000$; less 2\%:
USD~99\,470\,000.
\end{solution}

\begin{solution}{exo:m2:repo-and-specials:3}
\omterm{def:m2:repo-and-specials:special}{Specialness} $3.90 - 3.10 = 80$ basis points. Financing USD~100 million at
3.10\% instead of 3.90\% saves $100\,000\,000 \times 0.008/360 =
\text{USD}~2\,222.22$ a day (on the cash borrowed; scale by the dirty price for
the exact amount).
\end{solution}

\begin{solution}{exo:m2:repo-and-specials:4}
Coupon $100\,000\,000 \times 0.02125 \times 30/184 = 346\,467$; repo $100\,870\,911
\times 0.039 \times 30/360 = 327\,830$; carry USD~18\,637. DV01 USD~80\,444 per
basis point: breakeven 0.23 basis points. At 4.30\% the repo costs 361\,454 and
the carry is $-\text{USD}~14\,987$: the note now yields less than its financing,
and holding it is a bet that yields fall.
\end{solution}

\begin{solution}{exo:m2:repo-and-specials:5}
At 1\%: $\frac{1}{360} \times 0.01 \times 2 \times 100\,000\,000 = 5\,555.56$ a
day, USD~27\,777.78 for five days. At 3.75\%: $\max(3 - 3.75, 0) = 0$; nothing.
\end{solution}

\begin{solution}{exo:m2:repo-and-specials:6}
With $R = 0$: floor $-(3\% - 0) = -3\%$; shorts would pay up to 3\% a year to
borrow a scarce issue before failing. Today, with $R = 3.75\%$: floor 0.
Before 2009 there was no charge: failing cost only the forgone interest on
the proceeds, so no short would lend cash below 0\%, and when rates collapsed
at the end of 2008 that cost vanished too; fails became the cheapest
alternative to borrowing, and they multiplied.
\end{solution}

\begin{solution}{exo:m2:repo-and-specials:7}
After a fall of 1.5 points on the first day: the lender is owed $98\,853\,492$
plus one day's interest (10\,709), which requires collateral worth
$98\,864\,202/0.98$; the collateral is worth 99\,370\,911: call
USD~1\,510\,928. After 30 days at an unchanged price: the accrued interest
alone, grossed up by the haircut, $321\,274/0.98$: USD~327\,830. Real repo
books usually settle interest at maturity and margin on the principal; the
difference is a matter of each agreement's terms.
\end{solution}

\begin{solution}{exo:m2:repo-and-specials:8}
The Committee's decisions did not change on 17 September: the range stayed at
2.00--2.25\% and the administered rates with it. The market rate escaped the
range because reserves had become scarce for a day (tax payments and coupon
settlement drained 120 billion), so banks short of cash paid up and those with
cash lent it where rates were highest, in repo. The Fed responded by
\emph{adding} reserves through repo operations, and the next day it cut the
range by a quarter point: the opposite of tightening.
\end{solution}

\begin{solution}{pb:m2:repo-and-specials:1}
\textbf{1.} 3.45\%, 3.65\% and 3.80\%.
\textbf{2.} $100\,000\,000 \times 1.00870911 \times 0.0045 \times 30/360 =
\text{USD}~37\,827$.
\textbf{3.} $100\,870\,911 \times (30 \times 0.0045 + 30 \times 0.0025 + 30 \times
0.0010)/360 = \text{USD}~67\,247$.
\textbf{4.} The new issue is the market's hedging instrument, so shorts are
largest just after issue, while the supply that holders are willing to lend
is still building.
\textbf{5.} The shorts, through the cash they lend at the low rate to borrow the
note. They pay because failing to deliver would cost them their sale proceeds'
interest (and, when rates are low, the \omterm{def:m2:repo-and-specials:fails}{fails charge}), and because the note is
the hedge they need.
\textbf{6.} $V = 100.870911 \times 0.24/360 = 0.0672$ points per 100.
\textbf{7.} $0.0672 \times 32 = 2.15$ 32nds.
\textbf{8.} $0.0672/0.080444 = 0.84$ basis points.
\textbf{9.} About $4.210 - 0.008 = 4.202\%$.
\textbf{10.} At 4.195\% it is 0.7 basis points richer than financing alone
justifies: the rest is the price of liquidity.
\textbf{11.} Buy the off-the-run and sell the on-the-run in the same DV01; repo
the off-the-run at general collateral to finance it; reverse-repo the
on-the-run (lend cash, borrow the note) to deliver the short; hold until the
next auction makes the on-the-run old.
\textbf{12.} The long pays GC, 3.90\%, on its financing; the short earns only
the special rate on the cash it lends, 3.45\% then 3.65\% then 3.80\%: the
trade pays the \omterm{def:m2:repo-and-specials:special}{specialness}, which is what it is betting is priced too richly.
\textbf{13.} The short pays twice as much: 0.1345 points instead of 0.0672, and
the trade loses unless the yield spread moves in its favour by more.
\textbf{14.} A special rate of 0\% makes \omterm{def:m2:repo-and-specials:special}{specialness} 3.90\%: the short pays
USD~10\,928 a day on 100 million in forgone interest, and more if a squeeze
makes the note impossible to borrow at all.
\textbf{15.} Zero (\cref{prop:m2:repo-and-specials:floor}): with the reference
rate above 3\% the \omterm{def:m2:repo-and-specials:fails}{fails charge} is zero and failing costs no more than a
0\% loan.
\textbf{16.} Liquidity: the benchmark trades in the tightest markets, is
accepted everywhere as collateral and as a hedge, and investors pay for that
beyond its financing value.
\textbf{17.} If a borrower fails to return it, the lender loses its own ability
to deliver; the \omterm{def:m2:repo-and-specials:fails}{fails charge} is what the failing party would pay, and at
today's rates it is zero, so the protection is gone.
\textbf{18.} More supply of the same security to lend: \omterm{def:m2:repo-and-specials:special}{specialness} falls.
\textbf{19.} Named result: \emph{the value of the special} over the auction
cycle is 0.0672 points per 100 (2.15 32nds), 0.84 basis points of yield, USD~67\,247
on USD~100 million.
\textbf{20.} A security so much in demand for borrowing that lenders of cash
accept less than the general rate to get it.
\end{solution}

\begin{solution}{iq:m2:repo-and-specials:1}
A sale of securities with an agreement to buy them back: a loan of cash against
collateral, or of collateral against cash. Every long bond position is
financed in it and every short is borrowed in it; its rate is the traders'
cost of carry, and its \omterm{def:m2:repo-and-specials:special}{specialness} is the price of the securities they are
short.
\iqlookfor{both directions of the trade, and financing as the link to every
position.}
\end{solution}

\begin{solution}{iq:m2:repo-and-specials:2}
Its specific \omterm{def:m2:repo-and-specials:rate}{repo rate} is below general collateral: cash lenders accept a
lower rate to obtain that security, usually because it is in demand to cover
shorts or for delivery. The holder benefits: it borrows cash more cheaply
against it, or lends the security out; the shorts pay.
\iqlookfor{\omterm{def:m2:repo-and-specials:special}{specialness} as a transfer from shorts to holders.}
\end{solution}

\begin{solution}{iq:m2:repo-and-specials:3}
On 16 September 2019 tax payments and a coupon settlement drained about 120
billion of reserves, to 1.34 trillion; on the 17th Treasury repo printed at
5.25\% and the federal funds rate rose above its range. The New York Fed
lent 53 billion in an overnight repo operation, the Committee cut rates the
next day, and from October it bought bills to rebuild reserves; in 2021 it
created a standing repo facility.
\iqlookfor{reserve scarcity as the cause, and the three responses.}
\end{solution}

\begin{solution}{iq:m2:repo-and-specials:4}
Carry: the coupon accrued over the period less the repo interest on the
dirty price (plus a small pull to par). Breakeven: the yield rise that wipes it
out, carry divided by DV01; for a ten-year at 4.20\% financed at 3.90\% about a
quarter of a basis point a month.
\iqlookfor{the formula, and the conversion into basis points via DV01.}
\end{solution}

\begin{solution}{iq:m2:repo-and-specials:5}
Netting: cleared trades facing the central counterparty offset each other, so
a dealer's matched book of repo and \omterm{def:m2:repo-and-specials:rate}{reverse repo} shrinks to its net, saving
balance sheet and capital; the central counterparty also removes bilateral
credit exposure. The client gets access to dealers' balance sheet on better
terms, and the clearing mandate will require clearing for most repo anyway.
\iqlookfor{netting of balance sheet as the reason, not only credit risk.}
\end{solution}

\begin{solution}{iq:m2:repo-and-specials:6}
Yes, for a special. A short that must deliver compares lending cash at the
special rate $s$ to borrow the bond, earning $s$, with failing, earning nothing
and paying the \omterm{def:m2:repo-and-specials:fails}{fails charge} $(3\% - R)^+$. It borrows while $s > -(3\% -
R)^+$, so special rates can fall to $R - 3\%$ when the reference rate is below
3\%, and to zero otherwise. General collateral rates themselves go negative
only when policy rates are negative.
\iqlookfor{the arbitrage against failing, and the role of the \omterm{def:m2:repo-and-specials:fails}{fails charge}.}
\end{solution}
